% Generated from paper/source. Do not edit by hand.
\begin{abstract}
Clinical reasoning is sequential: evidence arrives with sign, absence, timing, and context, while a differential diagnosis should remain revisable until the evidence is sufficient. We investigate whether representing diagnostic hypotheses as an explicitly evolving complex state offers a useful inductive bias for this setting. Q-Neuro maps each observed finding to a low-rank evidence-conditioned operator, composes those operators in observed order, and measures diagnosis probabilities from squared complex amplitudes. The construction is quantum-inspired mathematics implemented on ordinary hardware; it is neither a quantum computer nor a physical model of cognition or disease.

We introduce NeuroWorld, a fully synthetic causal benchmark with chronology twins, factorial disease mechanisms, explicit missingness, unseen generator worlds, irreducible ambiguity, active evidence acquisition, omitted diseases, hidden syndromes, and counterfactual pairs. Across 26 registered studies---25 complete and one preserved failure---we compare Q-Neuro with unordered, recurrent, attention, state-space, associative, graph, real-operator, two-channel real, Hamiltonian-style, dissipative, attractor, tensor, and density-state controls. In a five-world severity sweep at 1,000 training cases, the complex operator exceeds the strongest tested two-channel real control by 0.054--0.063 absolute top-1 accuracy, with world-level 95\% intervals above zero. Critical ablations show that noncommutative order, phase-sensitive readout, and observed-negative evidence each contribute to this result.

The same program identifies decisive boundaries. A tuned GRU is substantially more sample-efficient in the source world; a real operator represents irreducible ambiguity far better; source-fitted temperature scaling worsens shifted calibration; hidden-syndrome and omitted-disease separation are not uniquely complex; phase-coded gradient optimization does not improve the frontier; and a velocity-based hard exit collapses to fixed two-step truncation rather than adaptive reasoning. Frozen-state probes recover simulator factors, but conventional recurrent controls recover them as well or better. The principal contribution is therefore not a claim of clinical or quantum advantage. It is a falsification-led, reproducible framework for testing hypothesis-state computation under causal shift, together with a bounded empirical finding: under the tested synthetic worlds, ordered complex dynamics improve robustness while paying measurable costs in calibration, compute, and ambiguity.

\textbf{Keywords:} complex-valued neural networks; diagnostic reasoning; causal simulation; distribution shift; noncommutative operators; uncertainty; reproducible machine learning

\end{abstract}
